% ===== 第五段：工具调用风险度量与可解释评估 =====
\subsubsection{Agent工具调用风险度量与可解释评估相关研究现状}

随着Agent系统部署规模的扩大，仅依赖攻防视角已不足以满足实际安全评估需求，
需要建立独立的、过程级的风险度量框架。

\textbf{行业标准与威胁建模框架。}
OWASP LLM Top 10\cite{owasp2025llmtop10}将过度代理权限（excessive agency）、
不安全的工具设计和不受控的执行行为列为Agent系统的核心安全风险，
指出工具调用链的权限边界管理是当前工业界Agent部署的首要安全挑战。
Cloud Security Alliance的MAESTRO框架\cite{huang2025maestro}将Agent威胁建模分解为七个协调层，
指出权限边界管理和执行可追溯性是工业界部署Agent的主要安全挑战，
并提供了面向企业级Agent系统的结构化风险评估方法论。
NIST AI风险管理框架（AI RMF）\cite{nist2023airmf}强调AI系统的风险评估需要从可测量、可解释的指标出发，
而非仅依赖黑盒评分，为建立基于证据的Agent安全评估体系提供了政策层面的支撑。

\textbf{工具调用评估的不足。}
在学术研究层面，现有工作大多以攻击成功率（ASR）或任务完成率（TCR）为核心评价指标，
缺少对工具调用链过程的精细度量。
Deng等\cite{deng2025agentsunderthreat}明确指出，
当前AI Agent安全基准缺乏统一的设计标准，
且ToolEmu等工具调用评测系统识别出的失效模式——
如Agent执行需要过度工具权限的操作、在未经用户确认情况下执行高风险命令——
说明工具调用链的安全性必须从权限边界、操作可逆性和用户确认等维度进行过程级评估，
而非仅从最终结果判断。
Chu等\cite{chu2026systematic}进一步指出，
所有已发布的Agent安全基准仅覆盖到T1（单次调用检测）或T2（有限序列检测）层面，
T3（跨步骤组合风险度量）和T4（过程级可解释评估）方面的进展在结构上不可测量，
即现有评估方法缺乏区分不同防御策略在组合风险检测上差异的能力。

\textbf{面向序列与轨迹的评估探索。}
SeqAR\cite{yang2025seqar}在NAACL 2025上提出基于序列自动生成的越狱评估框架，
将工具调用过程的时序特征纳入评测视角，
在Agent安全评测方向取得了一定进展，但仍局限于特定攻击类型的评估。
Yu等\cite{yu2025survey}在SIGKDD 2025上对可信LLM Agent的威胁与对策进行了系统综述，
梳理了从单次调用到多步任务链的安全评估挑战，
指出可信度评估需要同时考虑行为合规性、可解释性和责任归因三个维度。

\textbf{可解释性与归因度量。}
在可解释性方面，de Witt等\cite{dewitt2026openchallenges}指出，
现有评价指标大多是基于结果（outcome-based）的，
无法捕捉Agent间协调的质量或效率；
Dehghantanha等\cite{dehghantanha2026sokattacksurfaceagentic}在SoK中提出了
可从Agent执行轨迹直接计算的可攻击度量（UAR/PAR/RRS），
为基于轨迹的风险量化提供了初步思路，但尚未形成面向通用工具调用链的系统性度量框架。
Wang等\cite{wang2025protect}在IEEE/IFIP DSN 2025上从防御视角探讨了
如何通过可解释性机制增强Agent系统对提示注入的抵抗能力，
指出可解释性不仅有助于检测攻击，也是事后追责的必要条件。
Pearce等\cite{pearce2025asleep}在CACM 2025上通过大规模实验揭示了
AI代码生成工具在安全敏感代码上的系统性缺陷，
说明工具调用生成的代码同样需要过程级安全审计，
而非仅依赖最终代码的静态分析。

综上，工业界的威胁建模框架已明确提出过程级风险度量的需求，
学术界对跨步骤组合风险和可解释评估的研究尚处于起步阶段，
建立一套低成本、可复现、可解释的工具调用链风险指标体系具有重要的理论价值和实践意义。
